% App.F — FPGA Bitstream + SHA-256
% φ² + φ⁻² = 3 · Trinity S³AI
% Author: Dmitrii Vasilev (ORCID 0009-0008-4294-6159)

\chapter*{Appendix F: FPGA Bitstream Archive}

\begin{figure}[H]
\centering
\includegraphics[width=0.92\textwidth,keepaspectratio]{F-fpga-bitstream.jpg}
\caption*{Figure --- FPGA bitstream archive --- SHA-256 manifest.}
\end{figure}

\addcontentsline{toc}{chapter}{Appendix F: FPGA Bitstream Archive}
\label{app:fpga-bitstream}

\begin{quote}\itshape
``The best algorithms are those for which the obvious implementation is already
optimal---and which are hard to improve upon even in hindsight.''
--- Robert E.~Tarjan, \emph{Data Structures and Network Algorithms}, 1983.
\end{quote}

\section*{Silicon that survived contact with reality}

Bitstreams are the moment theory meets silicon. Everything in the preceding
chapters---the \(\varphi\)-numeric coprocessor, the Period-Locked Monitor at
92~MHz, the sub-watt power envelope---had to survive synthesis, place-and-route,
and a live JTAG download onto a QMTech XC7A100T board before any claim could
be called hardware-validated. This appendix archives the build artefacts:
the \texttt{.bit} file, its SHA-256 digest, the Vivado version, and the
timing-closure summary. It is the chain of custody between the HDL source
and the measurement results in Chapter~32. If you are reproducing the
experiment, this appendix tells you exactly which bitstream was used and how
to verify its integrity before flashing.

\section*{F.1 Overview}

This appendix documents the hardware artefacts produced during the Trinity S\textsuperscript{3}AI
FPGA validation campaign (2026-05-05). The bitstream implements the\emph{Period-Locked Monitor}
core (Chapter~\ref{ch:plm}) and the φ-numeric coprocessor primitives described in
Chapters~\ref{ch:28}--\ref{ch:34}.

\textbf{Anchor invariant:} \(\varphi^2 + \varphi^{-2} = 3\).

\section*{F.2 Target Device}

\begin{table}[h]
\centering
\resizebox{\linewidth}{!}{%
\begin{tabular}{ll}
\toprule
\textbf{Field} & \textbf{Value} \\
\midrule
Board & QMTech XC7A100T Starter Kit \\
Actual device & XC7A200T (IDCODE \texttt{0x13631093}) \\
Manufacturer & Xilinx (JEDEC \texttt{0x049}) \\
Version & Silicon revision 1 \\
Package & FGG484 \\
Speedgrade & --2 \\
\bottomrule
\end{tabular}
}
\caption{FPGA device identification. Board is marked XC7A100T but JTAG IDCODE
corresponds to XC7A200T v1. The design (83~LUT, 27~FF) fits either device;\, no
resynthesis required.}
\end{table}

\textbf{IDCODE parsing} (IEEE~1149.1):
\begin{itemize}
  \item Bits 0: \texttt{1} (required by standard)
  \item Bits 1--11: \texttt{0x049} — Xilinx manufacturer
  \item Bits 12--27: \texttt{0x3631} — XC7A200T part
  \item Bits 28--31: \texttt{0x1} — silicon revision 1
\end{itemize}

\section*{F.3 Bitstream Integrity}

\begin{table}[h]
\centering
\resizebox{\linewidth}{!}{%
\begin{tabular}{ll}
\toprule
\textbf{Field} & \textbf{Value} \\
\midrule
Filename & \texttt{design.bit} \\
SHA-256 & \texttt{8536e265...d77352b} \\
Tool & Vivado 2023.2 \\
Strategy & Default (Flow\_PerfOptimized\_high) \\
LUT utilization & 83 / 128,000 (0.06\%) \\
FF utilization & 27 / 256,000 (0.01\%) \\
BRAM & 0 \\
DSP & 0 \\
\bottomrule
\end{tabular}
}
\caption{Bitstream file metadata. Full SHA-256 stored in
\filepath{trinity-fpga/bitstream/design.bit.sha256}.}
\end{table}

\section*{F.4 Configuration Status Register (STAT)}

After successful programming via ESP32 XVC WiFi bridge
(Section~\ref{sec:xvc-bridge}), the STAT register was read using
\texttt{openFPGALoader --read-register STAT}:

\begin{verbatim}
Raw STAT register: 0x401079FC

CRC Error     : No CRC error       ✓
Part Secured  : 0                  ✓
MMCM Lock     : 1                  ✓
DCI Match     : 1                  ✓
EOS           : 1 (End of Startup) ✓
GWE           : 1 (Global Wr En)   ✓
GHIGH_B       : 1                  ✓
Done          : 1                  ✓
DEC Error     : 0                  ✓
ID Error      : 0                  ✓
\end{verbatim}

All flags nominal. The device entered operational state without errors,
confirming bitstream integrity and correct CRC check by the FPGA configuration
logic.

\section*{F.5 Programming Method}\label{sec:xvc-bridge}

The bitstream was loaded via a custom ESP32-based Xilinx Virtual Cable (XVC)
server operating over 802.11 WiFi (IP~\texttt{192.168.1.30}, port~2542).
Full details in Appendix~J and Chapter~\ref{ch:33}.

\begin{verbatim}
openFPGALoader --cable xvc-client \
  --ip 192.168.1.30 --port 2542 \
  --bitstream design.bit
\end{verbatim}

\section*{F.6 Reproducibility}

The complete firmware, XDC constraints, and bitstream are archived in
\filepath{gHashTag/trinity-fpga} (Apache~2.0). To reproduce:

\begin{enumerate}
  \item Flash ESP32 with \filepath{firmware/xvc-esp32/xvc-esp32.ino}
  \item Connect 5-wire JTAG harness per Appendix~I
  \item Run \texttt{openFPGALoader} command above
  \item Verify STAT\,=\,\texttt{0x401079FC}
\end{enumerate}
